%% ===== Introduction (~500-650 words) =====
\section{Introduction\cb{target ~570 words}\wc{intro}}\label{sec:intro}
% NOTE: the TorchCell overview schematic (fig:overview) is parked in the
% Supplementary information for now (see sections/backmatter.tex); the Intro and
% Results R1 reference it as a supplementary figure.
Mapping genotype to phenotype is a central goal of biology and the rate-limiting step of
metabolic engineering, where the targets of interest---titre, rate, and yield---are
polygenic phenotypes that emerge from many genes acting across transcription, translation,
and metabolism~\cite{stephanopoulosMetabolicEngineeringPrinciples1998,bordbarConstraintbasedModelsPredict2014}. Unlike monogenic disorders, such
phenotypes cannot be read from a single gene: their causal structure is obscured by
epistasis and by the cellular layers that separate genotype from
measurement~\cite{costanzoGlobalGeneticInteraction2016a,kuzminSystematicAnalysisComplex2018a}. Growth is the most tractable of
them---assays are cheap and scalable, gene--gene interactions are directly quantifiable, and
the same readout spans applications from suppressing proliferation in disease to maximizing
growth-coupled production.

Deep learning has transformed protein structure and function prediction, yet systems biology
and metabolic engineering lag~\cite{sapovalCurrentProgressOpen2022}. Two obstacles recur. First, phenotype
data are scattered across small, heterogeneous studies with no shared schema, and broad
genome-wide screens rarely interoperate with the deep, pathway-focused campaigns of metabolic
engineering~\cite{volkMetabolicEngineeringMethodologies2022,presnellSystemsMetabolicEngineering2019}. Second, although multimodal measurements
improve prediction, combining more than a few modalities is difficult and rarely
scales~\cite{culleyMechanismawareMultiomicMachinelearning2020}. Recent foundation cell models pretrain on single-cell
gene-expression matrices and fine-tune for phenotype~\cite{cuiScGPTBuildingFoundation2024,kalfonScPRINTPretraining50MCells2025},
but industrial microbes lack such data; what they do have is diverse perturbation data and
curated interaction graphs, including metabolism.

Existing approaches leave this gap open: constraint-based flux analysis captures metabolism
but misses epistasis~\cite{bordbarConstraintbasedModelsPredict2014,merzbacherAccuratePredictionGene2025}, and recent perturbation
transformers do not consistently beat linear baselines~\cite{ahlmann-eltzeDeeplearningbasedGenePerturbation2025}. Our own
classical baselines make the point sharply---with enough data they predict knockout fitness
from almost any gene representation, yet fail on gene interactions (Supplementary Fig.~\ref{fig:classical-ml})---so
fitness alone does not justify a deep model, but the harder, structured labels do.

Here we present TorchCell, a schema-grounded knowledge graph that curates heterogeneous
\textit{S.\ cerevisiae} phenotype data into queryable multimodal datasets, in the
data-standardization spirit of TorchGeo~\cite{stewartTorchGeoDeepLearning2022a}, together with the Cell
Graph Transformer (CGT): a model that encodes a cell once, applies perturbations as a learned
operator in embedding space, and folds biological networks into the training objective as a
soft prior rather than a hard graph. It predicts fitness, gene interactions, gene expression,
and morphology jointly, and needs only a genome---not an interaction graph---to run. We
outline the data model and architecture in Supplementary Fig.~\ref{fig:overview}.
